\subsection{RQ2: Contribution of Admission and Pacing}
\label{sec:rq2}

% Notes: docs/exhibits.md#tab_rq2_ablation
\begin{table}[t]
  \centering
  \setlength{\abovecaptionskip}{2pt}%
  \setlength{\belowcaptionskip}{3pt}%
  \caption{Contributions of admission and pacing (RQ2).}
  \label{tab:rq2_ablation}

  {\scriptsize
  \renewcommand{\arraystretch}{1.05}%
  % Shaded rows: fill the gap before each rule with the row color, and
  % keep the shading inside the outer edges of the first and last columns.
  \makeatletter
  \newcommand{\rqablgap}{\ifx\CT@row@color\relax\hspace{1.5pt}%
    \else{\let\CT@color\color\CT@row@color\vrule\@width1.5pt}\fi}%
  \newcommand{\rqablclipL}{\gdef\CT@cell@color{\@tempdimb\z@
    \global\let\CT@cell@color\relax}}%
  \newcommand{\rqablclipR}{\gdef\CT@cell@color{\@tempdimc\z@
    \global\let\CT@cell@color\relax}}%
  \makeatother
  \newcolumntype{V}{!{\rqablgap\vrule}}%
  \newsavebox{\rqablbox}%
  \newcommand{\rqabltabular}{%
  \begin{tabular}{@{}%
      >{\rqablclipL\raggedright\arraybackslash}p{46pt}V%
      >{\raggedleft\arraybackslash}p{29pt}%
      >{\raggedleft\arraybackslash}p{29pt}V%
      >{\raggedleft\arraybackslash}p{23pt}%
      >{\raggedleft\arraybackslash}p{23pt}V%
      >{\raggedleft\arraybackslash}p{34pt}%
      >{\rqablclipR\raggedleft\arraybackslash}p{33pt}%
    @{}}
    \toprule
    \multirow[c]{2}{=}[-4.5pt]{\centering\makecell[c]{\textbf{Control}\\\textbf{configuration}}}
      & \multicolumn{2}{cV}{\makecell[c]{\textbf{Deadline safety}}}
      & \multicolumn{2}{cV}{\makecell[c]{\textbf{Stage execution}}}
      & \multicolumn{2}{c}{\makecell[c]{\textbf{Pacing cost}}} \\
    \cmidrule(lr){2-3}\cmidrule(lr){4-5}\cmidrule(l){6-7}
      & \multicolumn{1}{c}{\makecell[c]{\textbf{Survived}\\\textbf{runs}}}
      & \multicolumn{1}{cV}{\makecell[c]{\textbf{Captures}\\\textbf{(\%)}}}
      & \multicolumn{1}{c}{\makecell[c]{\mbox{\boldmath$M$}\\\textbf{(\%)}}}
      & \multicolumn{1}{cV}{\makecell[c]{\mbox{\boldmath$S$}\\\textbf{(\%)}}}
      & \multicolumn{1}{c}{\makecell[c]{\textbf{Activated}\\\textbf{(\%)}}}
      & \multicolumn{1}{c@{}}{\makecell[c]{\mbox{\boldmath$\mathit{delay}$}\textbf{ P50}\\\textbf{(ms)}}} \\
    \midrule

    \multicolumn{7}{@{}l}{\textbf{S26 Ultra, 12MP normal, Lv4}} \\
    \addlinespace[0pt]
    No control     & 0/10   & 29.0  & 29.0 & 29.0 & 0 & -- \\
    Admission only & 10/10  & 100.0 & 29.3 & 33.7 & 0 & -- \\
    Pacing only    & 0/10   & 44.3  & 44.3 & 44.3 & 51.1 & 242 \\
    \rowcolor{gray!20}
    \textbf{Full}  & \textbf{10/10} & \textbf{100.0}
                   & \textbf{58.0} & \textbf{95.7} & \textbf{37.9} & \textbf{437} \\
    \midrule

    \multicolumn{7}{@{}l}{\textbf{S26 Ultra, 24MP memory, Lv4}} \\
    \addlinespace[0pt]
    No control     & 0/10   & 24.3  & 24.3 & 24.3 & 0 & -- \\
    Admission only & 8/10   & 87.3  & 24.3 & 51.0 & 0 & -- \\
    Pacing only    & 0/10   & 28.1  & 28.1 & 28.1 & 58.4 & 483 \\
    \rowcolor{gray!20}
    \textbf{Full}  & \textbf{10/10} & \textbf{100.0}
                   & \textbf{44.5} & \textbf{77.0} & \textbf{36.7} & \textbf{400} \\
    \midrule

    \multicolumn{7}{@{}l}{\textbf{S26, 12MP normal, Lv4}} \\
    \addlinespace[0pt]
    No control     & 0/10   & 42.3  & 42.3 & 42.3 & 0 & -- \\
    Admission only & 7/10   & 90.7  & 44.3 & 51.3 & 0 & -- \\
    Pacing only    & 1/10   & 61.4  & 61.4 & 61.4 & 44.7 & 387 \\
    \rowcolor{gray!20}
    \textbf{Full}  & \textbf{10/10} & \textbf{100.0}
                   & \textbf{56.0} & \textbf{74.7} & \textbf{53.1} & \textbf{297} \\
    \bottomrule
  \end{tabular}}%
  \fittabcolsep{\rqablbox}{\rqabltabular}{\columnwidth}{12}%
  \resizebox{0.95\columnwidth}{!}{\rqabltabular}
  }
\end{table}


Table~\ref{tab:rq2_ablation} compares admission and pacing separately and together in three cells starting at the static thermal guard's level-4 cutoff.
\emph{No control} enables neither component, the same configuration as Baseline in Table~\ref{tab:rq1_controller_behavior}.
\emph{Admission only} and \emph{Pacing only} enable only the named component, and \emph{Full} enables both.
The Captures column reports the percentage of the 30 planned captures per run completed on time.
The \(M\) and \(S\) columns report the percentages completed on time with the corresponding stage executed, using the same denominator.

\myparagraph{Limits of single-component control}
Only \emph{Full} completed all planned captures on time in all three cells.
On-time completion rates were \mbox{24.3--42.3\%} for \emph{No control} and \mbox{28.1--61.4\%} for \emph{Pacing only}.
\emph{Admission only} reached 100\% in the S26 Ultra 12MP cell but 87.3--90.7\% in the other two cells.

In the S26 Ultra 12MP cell, \emph{Admission only} met every deadline while skipping each optional stage in most captures.
Its \(M\) and \(S\) execution rates were 29.3\% and 33.7\%, compared with 58.0\% and 95.7\% under \emph{Full}.
In the S26 Ultra 24MP and S26 12MP cells, \emph{Admission only} also executed both stages less often than \emph{Full}, yet Capture Timeout still occurred.

\myparagraph{Coordinated control}
\emph{Full} executed both \(M\) and \(S\) more often than \emph{Admission only} in every cell.
However, this retention came with a responsiveness cost.
\emph{Full} activated pacing on 36.7--53.1\% of transitions, with a median positive applied delay of 297--437~ms.
These medians were 37--45\% of the corresponding shot-to-shot intervals with parallel capture disabled.
